% Generated draft: requires mathematical and source review.
\documentclass[11pt]{article}
\usepackage[margin=25mm]{geometry}
\usepackage{fontspec}
\setmainfont{Latin Modern Roman}
\usepackage{amsmath,amssymb,mathtools,unicode-math}
\setmathfont{Latin Modern Math}
\usepackage{enumitem,needspace,xurl,hyperref}
\setlength{\parindent}{0pt}
\setlength{\parskip}{5pt}
\setlength{\emergencystretch}{3em}
\newcommand{\N}{\mathbb{N}}
\newcommand{\Z}{\mathbb{Z}}
\newcommand{\Q}{\mathbb{Q}}
\newcommand{\R}{\mathbb{R}}
\newcommand{\C}{\mathbb{C}}
\newcommand{\F}{\mathbb{F}}
\newcommand{\A}{\mathbb{A}}
\newcommand{\PP}{\mathbb{P}}
\newcommand{\E}{\mathbb{E}}
\newcommand{\eps}{\varepsilon}
\newcommand{\dd}{\,\mathrm{d}}
\begin{document}
\section*{UnsolvedMath ID: 30001042}
\textbf{Source problem code:} catalogue entry
\subsection{Definitions and mathematical statement}
For metric spaces $(Y,d_Y)$ and $(Z,d_Z)$, a bi-Lipschitz bijection is a bijection $f:Y\to Z$ such that for some $C\geq1$,
\[
C^{-1}d_Y(x,y)\leq d_Z(f(x),f(y))\leq C d_Y(x,y)\quad(x,y\in Y).
\]
For finitely generated groups, use word metrics from finite generating sets; their bi-Lipschitz class does not depend on the finite generating set. Equip a product with $\Z^2$ with the corresponding product word metric. No map is required to be a homomorphism and the constant may depend on the groups.

The source bundles four questions:
\begin{enumerate}
\item For every finitely generated group $G$, is there a bi-Lipschitz bijection $G\to G\times\Z^2$?
\item For every finite-index subgroup $G_1\leq G$ of a finitely generated group, is there a bi-Lipschitz bijection $G_1\to G$?
\item For two uniform lattices $L_1,L_2$ in the same Lie group, is there a bi-Lipschitz bijection $L_1\to L_2$? A uniform lattice is a discrete cocompact subgroup. Word metrics apply when the lattices are finitely generated; the source does not give the Lie-group hypotheses or an alternate metric convention otherwise.
\item If a compact three-manifold $M$ admits a partially hyperbolic diffeomorphism with invariant splitting
\[
TM=E^s\oplus E^c\oplus E^u,\qquad
\dim E^s=\dim E^c=\dim E^u=1,
\]
what restrictions follow for the topology of $M$? Conventionally, partial hyperbolicity means uniform contraction on $E^s$, uniform expansion on $E^u$, and domination of the center by the unstable and of the stable by the center, after suitable iterates and choice of metric.
\end{enumerate}
The first literal quantifier includes finite groups, for which there cannot be a bijection to $G\times\Z^2$ on cardinality grounds. The second likewise needs finite cases checked if $G_1$ is proper. These are formulation obstructions to treating the supplied unrestricted wording as a verified open universal assertion. No replacement stabilization, exclusion of groups, or extra hypothesis is guessed. Quasi-isometry is not substituted for bi-Lipschitz bijectivity.

The final topology question is a separate broad task, not a consequence of the preceding group questions. The source does not specify closed versus boundary-bearing manifolds, orientability, diffeomorphism regularity, precise domination convention or the type of topological characterization sought.

The supplied triage checked 2026-08-21 labels the bundle open and says the stabilization, finite-index and broad topology questions remain unknown in recent literature. That claim is retained as attributed curation, with the literal finite-case issue explicitly flagged rather than accepted as a mathematical verification.
\subsection{Short English statement}
Preserve three proposed bi-Lipschitz group comparisons and the separate topology question for three-dimensional partially hyperbolic dynamics. The universal group wording has finite-case obstructions, and the lattice metrics and dynamical conventions need clarification.
\subsection{Sources}
\begin{itemize}
\item Local UnsolvedMath registry, permanent ID 30001042, source code catalogue entry; supplied complete record and corrected-verified notation-cleanup note. Literature checked 2026-08-21: the assessment is supplied dated curation, not independently verified here.
\item Geometric Group Theory, Hyperbolic Dynamics and Symplectic Geometry (2009), Oberwolfach Reports, DOI 10.4171/owr/2008/42, pp. 2337--2406: \url{https://doi.org/10.4171/owr/2008/42}. Proposed-year metadata gives 2008.
\item Petrunin, metric geometry monograph: \url{https://amathr.org/wp-content/uploads/2025/01/Petrunin_AMRResMonograghs2Ed2Pr.pdf}.
\end{itemize}
\end{document}
